\chapter{数值算法设计与收敛判据}

\section{算法总体流程}
本文数值流程分为四个阶段：
\begin{enumerate}
  \item 读取测角和测距观测数据；
  \item 由三测站测距几何定位得到初值 $\bm p_0$；
  \item 采用牛顿法或法方程迭代更新估计；
  \item 输出最优解、目标函数值、矩阵及精度指标。
\end{enumerate}

\section{牛顿法}
\subsection{迭代公式}
在第 $k$ 步，计算
\[
\bm g_k=\nabla F(\bm p_k),\quad
\bm H_k=\nabla^2F(\bm p_k),
\]
解线性系统
\[
\bm H_k\Delta\bm p_k=\bm g_k,
\]
并更新
\[
\bm p_{k+1}=\bm p_k-\Delta\bm p_k.
\]

\subsection{收敛判据}
定义相对步长
\begin{equation}
\eta_k=\frac{\|\Delta\bm p_k\|_2}{\|\bm p_{k+1}\|_2}.
\end{equation}
当 $\eta_k<10^{-10}$ 时判定收敛。若迭代超过上限（代码中设置额外保护），则返回失败状态。

\subsection{优缺点}
\begin{itemize}
  \item 优点：局部二次收敛，收敛后可直接得到海森矩阵；
  \item 缺点：对初值与海森矩阵正定性敏感；
  \item 适用：观测质量较高、初值较好、维度中小的场景。
\end{itemize}

\section{法方程迭代法（代码中的 LM 方法）}
\subsection{线性化思路}
在 $\bm p_k$ 处线性化残差函数：
\[
r_i(\bm p_k+\Delta\bm p)\approx r_i(\bm p_k)-\bm J_i\Delta\bm p.
\]
令加权残差最小，可得法方程
\begin{equation}
\bm A_k\Delta\bm p_k=\bm B_k,
\label{eq:normal_eq}
\end{equation}
其中
\[
\bm A_k=\sum_i w_i\bm J_i^{\T}\bm J_i,\quad
\bm B_k=\sum_i w_i r_i\bm J_i^{\T}.
\]
更新规则为
\[
\bm p_{k+1}=\bm p_k+\Delta\bm p_k.
\]

\subsection{与高斯牛顿关系}
式 \eqref{eq:normal_eq} 本质上是高斯牛顿框架。代码命名为 LM 方法，但未显式加入阻尼参数 $\lambda$。若后续需要增强鲁棒性，可拓展为
\[
(\bm A_k+\lambda_k\bm I)\Delta\bm p_k=\bm B_k.
\]

\subsection{优缺点}
\begin{itemize}
  \item 优点：避免显式二阶导数，计算结构清晰；
  \item 缺点：强非线性区域可能收敛变慢；
  \item 适用：导数信息可得且希望快速实现工程求解的场景。
\end{itemize}

\section{算法伪代码}
\begin{algorithm}[H]
\caption{融合测角测距的牛顿迭代定位}
\begin{algorithmic}[1]
\State 输入观测集 $S_1,S_2$，初值 $\bm p_0$
\For{$k=0,1,\ldots$}
\State 计算 $\bm g_k=\nabla F(\bm p_k)$ 和 $\bm H_k=\nabla^2F(\bm p_k)$
\State 求解 $\bm H_k\Delta\bm p_k=\bm g_k$
\State $\bm p_{k+1}\gets \bm p_k-\Delta\bm p_k$
\If{$\|\Delta\bm p_k\|_2/\|\bm p_{k+1}\|_2<10^{-10}$}
\State 返回 $\bm p_{k+1}$
\EndIf
\EndFor
\State 返回失败状态
\end{algorithmic}
\end{algorithm}

\begin{algorithm}[H]
\caption{法方程迭代定位}
\begin{algorithmic}[1]
\State 输入观测集 $S_1,S_2$，初值 $\bm p_0$
\For{$k=0,1,\ldots$}
\State 线性化残差，累加构造 $\bm A_k,\bm B_k$
\State 求解 $\bm A_k\Delta\bm p_k=\bm B_k$
\State $\bm p_{k+1}\gets \bm p_k+\Delta\bm p_k$
\If{$\|\Delta\bm p_k\|_2/\|\bm p_{k+1}\|_2<10^{-10}$}
\State 计算 $\bm A_k^{-1}$ 并返回
\EndIf
\EndFor
\State 返回失败状态
\end{algorithmic}
\end{algorithm}

\section{复杂度分析}
设总观测数 $N=2m+n$（测角每站贡献两个残差）。每次迭代主要代价在残差与导数累计，复杂度约为 $\mathcal O(N)$；线性求解为 $3\times3$ 小系统，复杂度常数量级。因此算法总耗时主要由迭代次数决定。

\section{停止准则与异常处理}
为提高工程可用性，建议采用以下多重停止准则：
\begin{enumerate}
  \item 参数增量准则：$\|\Delta\bm p\|/\|\bm p\|<\varepsilon_1$；
  \item 目标下降准则：$|F_{k+1}-F_k|<\varepsilon_2$；
  \item 梯度准则：$\|\nabla F\|<\varepsilon_3$；
  \item 最大迭代次数准则。
\end{enumerate}
此外可在每次迭代监测矩阵行列式与条件数，必要时触发正则化或重启策略。

\section{本章小结}
本章完成了两类求解器的算法化表述，并给出各自优缺点与实现要点。后续实验将从收敛结果、精度表现和鲁棒性等方面进行对比分析。
